% This makes the macros `\inputleanmodule` and `\inputleannode` available.
\input{../../.lake/build/blueprint/library/Cslib}

\chapter{Alpha equivalence equalities}

\inputleanmodule{Cslib.Languages.LambdaCalculus.Named.Untyped.AlphaEquivVariants.Defs}
\inputleanmodule{Cslib.Languages.LambdaCalculus.Named.Untyped.Basic}
\inputleanmodule{Cslib.Languages.LambdaCalculus.Named.Untyped.AlphaEquivProperties}
\inputleanmodule{Cslib.Languages.LambdaCalculus.Named.Untyped.AlphaEquivVariants.Swap}
\inputleanmodule{Cslib.Languages.LambdaCalculus.Named.Untyped.AlphaEquivVariants.SwapSubst}
\inputleanmodule{Cslib.Languages.LambdaCalculus.Named.Untyped.AlphaEquivVariants.Equalities}

\chapter{De Bruijn terms really do work}

This chapter follows M.~Norrish and R.~Vestergaard, \emph{Proof Pearl: de Bruijn Terms Really Do
Work} (TPHOLs 2007). Results are numbered as in that paper.

\section{Named abstraction (Section 2.1)}

\inputleannode{def:nv-lambda-alpha}
\inputleannode{def:nv-lambda-alpha-subst}
\inputleannode{def:nv-beta-alpha}

\section{De Bruijn's anonymous abstraction (Section 2.2)}

\inputleannode{def:nv-db}
\inputleannode{def:nv-lift}
\inputleannode{def:nv-nsub}
\inputleannode{def:nv-1}

\section{An intermediate formalism (Section 4.1)}

\inputleannode{def:nv-4}
\inputleannode{def:nv-dfv}
\inputleannode{def:nv-5}
\inputleannode{lem:nv-6-aux}
\inputleannode{lem:nv-6}
\inputleannode{def:nv-7}

\section{Nominal reasoning with indices (Section 4.2)}

\inputleannode{def:nv-8-inc}
\inputleannode{def:nv-8}
\inputleannode{lem:nv-supp}
\inputleannode{lem:nv-9}
\inputleannode{thm:nv-10-sub}
\inputleannode{thm:nv-10-dlam}
\inputleannode{thm:nv-10}
\inputleannode{lem:nv-dfv-dlam}
\inputleannode{lem:nv-sub-dlam}
\inputleannode{def:nv-f-raw}
\inputleannode{lem:nv-f-respects-alpha}
\inputleannode{def:nv-f}
\inputleannode{lem:nv-f-injective}
\inputleannode{lem:nv-f-surjective}
\inputleannode{thm:nv-11}
\inputleannode{lem:nv-12-subst}
\inputleannode{thm:nv-12}

\section{The homomorphism between dB and dB' (Section 4.3)}

\inputleannode{lem:nv-13}
\inputleannode{def:nv-lifting-pm}
\inputleannode{lem:nv-14}
\inputleannode{lem:nv-sub-fresh-perm}
\inputleannode{lem:nv-dlam-inj}
\inputleannode{thm:nv-15}
\inputleannode{lem:nv-16}
\inputleannode{lem:nv-17}
\inputleannode{lem:nv-beta-equivariant}
\inputleannode{lem:nv-18}
\inputleannode{thm:nv-19}
\inputleannode{thm:nv-20-hom}
\inputleannode{thm:nv-20}

\section{Eta-reduction (Section 4.4)}

\inputleannode{def:nv-eta-alpha}
\inputleannode{def:nv-neta}
\inputleannode{def:nv-eta}
\inputleannode{lem:nv-nsub-lift}
\inputleannode{lem:nv-lift-nsub}
\inputleannode{thm:nv-eta-coincide}
\inputleannode{thm:nv-eta-iso}

\section{Using the isomorphism}

\inputleannode{cor:nv-church-rosser}
\inputleannode{cor:nv-normal-form-unique}
